\section{Structural Estimation of DDCs\label{sec:ddc}}

The development of new recursive methods for solving DPs and MDPs in the 1950s and 60s sparked
a revolution in economics because these tools enabled economists to model sequential
decision making under uncertainty, dynamic strategic interactions, and dynamic equilibria 
in markets which are key features of nearly all economic problems.  Starting in the 1960s
DP facilitated the development of new dynamic economic theories that showed 
how expectations of the future affect the current choices of rational forward looking agents, 
resulting in new dynamic economic theories embodying {\it rational expectations\/} \citet{Muth:1961}. 
Since this new framework had the flexibility to predict a wide range of behavior, the natural
next step was to develop empirical methods for estimating and testing how well these new theories
can explain data. This lead to a new area of econometrics that can be called {\it dynamic structural estimation\/} (DSE)
which differed from existing econometric methods that could be classified as {\it reduced form estimation\/} (RFE). 
The term ``structural'' refers to the attempt to infer the underlying preferences and beliefs of economic
agents rather than simply summarizing their behavior, the main goal of reduced form econometrics.
In terms of the MDP framework, we can consider an agent's decision rule $\pi(a|s)$ to be the ``reduced form model''
because it captures their behavior, whereas structural econometrics tries to infer the underlying
{\it structure\/} i.e. the preferences and beliefs $\{\beta,r,p\}$ as well as the implied optimal decision rule $\pi$.\footnote{\footnotesize
The earliest development of structural econometrics dates back to work at the Cowles Foundation in thee 1940s where
econometric methods were developed for inference on {\it linear simultaneous equation systems\/} which are
static models of market equilibrium. In this work  the structure was the coefficients of linear market
supply and demand curves, whereas the reduced form were linear regressions that predict how equilibrium
prices and quantities $(p,q)$ shift with various variables $x$ that shift the supply and demand curves. 
This was the earliest work showing how it could be possible to infer the structure (i.e. the coefficients of
the supply and demand curves) using only data $(p,q,x)$ on the intersections of the supply and demand curves in 
different market states $x$.}

\citet{lucas:1976} provide a practical rationale for why we do structural econometrics. 
He  argued that reduced-form econometric models produce misleading predictions of the effect of
economic policy changes, whereas structural models that estimate the underlying preferences of consumers
and the production technologies of firms are more likely to accurately predict the effects of 
policy changes. Policy evaluation requires counterfactual predictions of how 
economic agents and the economy will react to changes in economic policy and the economic environment.
To explain Lucas's argument in the MDP framework, let $\rho$ devote a set of environment
or policy variables that affect the preferences (or profits) and beliefs of economic agents. We can
write the reward function as $r(s,a,\rho)$ and beliefs as $p(s'|s,a,\rho)$ to emphasize the dependence
of these quantities on $\rho$. We can solve the MDP model for different values of $\rho$ to 
to make counterfactual predictions of how behavior $\pi(a|s,\rho)$ and welfare $V(s,\rho)$ change in response
to changes in $\rho$. As noted, when there is insufficient variation in $\rho$ (or we want to consider changes in $\rho$
beyond the range of historical experience), reduced form methods cannot provide good estimates of how $\pi(a|s,\rho)$
as $\rho$ changes, nor can they predict how it affects welfare. However it is sometimes possible to know how $\rho$ affects $r(s,a,\rho)$ 
and $p(s'|s,a,\rho)$ even when there has been little or no variation in $\rho$ in the past. The introduction discussed the 
example of new car taxation in Denmark discussed in the introduction where $\rho$ is the tax rate on new cars. 
Under the assumption that  $r(s,a,\rho)$ additively separable in $(s,a)$ and
$\rho$ \cite{GIMRS:2022} used structural estimation to recover the component of $r(s,a,\rho)$ that depends
only on $(s,a)$, and with this, they could calculate  $r(s,a,\rho)$ for any other possible tax rate $\rho$, due to the
their assumption about its functional form. In contrast, we rarely have a good prior
known of how $\rho$ affects the decision rule  $\pi(a|s,\rho)$ 
but structural estimation of $r(s,a,\rho)$ and $p(s'|s,a,\rho)$ enables us to calculate $\pi(a|s,\rho)$  for any policy $\rho$
by solving the MDP. These solutions constitute  counterfactual predictions of how car holding and trading behavior and welfare
changes in response to changes in the new car tax rate $\rho$.\footnote{\footnotesize Since changes in tax rates $\rho$ aslo affect prices
of new and used cars, a counterfactual prediction also has to account for how $\rho$ affects the entire equilibrium
in the Danish car market, so both new and used car prices are implicit functions of $\rho$. These {\it equilibrium
counterfactuals\/} are also captured by the structural model of \citet{GIMRS:2022}.}

The earliest work on estimation of dynamic structural models includes 
\citet{sargent:1978}, who estimated the dynamic linear quadratic model\footnote{\footnotesize That is, the law of motion for $s_t$ is a linear autoregression 
and $r$ is a quadratic function of $(s_t,a_t)$.} of firms' demand for workers by maximum likelihood. 
In the remainder of this section we provide a selective summary of the literature on structural estimation of DDCs that has the closest connection to 
the literature on IRL. Key references include \citet{rust:1987}, \citet{rust:1988}, \citet{rust:1994}, \citet{AguirregabiriaMira:2002},
and the survey by \citet{AguirregabiriaMira:2010}. 

The goal is to infer the structural objects
$\{\beta,r,p\}$ from data on the trajectories of choices and states on a random sample of individuals 
using a modified version of the finite state and action MDP discussed in section 2.  
The modification is motivated by the {\it statistical degeneracy\/} of the optimal decision rule $\pi(a|s)$ 
in equation (\ref{eq:bellman_policy}), i.e. it is generically a pure strategy (i.e. deterministic function of $s$). 
For most choices we actually observe, different individuals will choose different actions in the same
observed state $s$.  In order to explain this heterogeneity in behavior, we consider a modified MDP with an augmented
state variable $x=(s,\epsilon)$ where the individual
observes both $s$ and $\varepsilon$ (and thus uses a pure strategy) but we only observe $s$  but not $\varepsilon$, which 
is a vector of idiosyncratic preference shocks that we do not observe. As a result, from our standpoint
individuals appear to be using a mixed strategy. If there is sufficient variation
in the unobserved $\varepsilon$ component, we can derive a {\it conditional choice probability\/} (CCP) $\pi(a|s)$  from the overall (pure)
strategy $\pi(a|s,\varepsilon)$ that admits positive probabilities for all feasible choices $a$.

Building on the contributions of \citet{mcfadden:1973} for static discrete choice
models, we can derive an especially simple expression for the CCP $\pi(a|s)$ given by
the {\it multinomial logit model\/} or ``softmax'' function in the RL literature.
Suppose the choice sets depend only on $s$ and are finite, and assume that for each
$s$, $\varepsilon$ is a vector with the same length as $|A(s)|$ the number of actions, so for
$\varepsilon(a)$ is the component of $\varepsilon$ associated with action $a \in A(s)$. Assume that 
the reward for action $a$ in state $a$ $(s,\varepsilon)$ can be written as $r(s,a)+\varepsilon(a)$. We call
this the Additive Separability (AS) assumption. We also assume
 that the transition density $p(s',\varepsilon'|s,\varepsilon,a)$ factors as $g(\varepsilon'|s')p(s'|s,d)$,
for a conditional density $g(\varepsilon'|s')$.
We call this the Conditional Independence (CI) assumption. Finally, if 
we assume $\varepsilon$ has a multivariate Gumbel or Type 1 extreme value distribution given by
\begin{equation}
    F(\varepsilon|s) = \prod_{a \in A(s)} \exp\{-\exp\{-(\varepsilon(a)-\mu)/\sigma\}\},
\label{eq:evd}
\end{equation}
where we set $\mu=-\sigma\gamma$ where $\gamma$ is Euler's constant, then $E\{\varepsilon(a)|s\}=0$ for $a \in A(s)$. 
We call this the Extreme Value (EV) assumption. Let $V(s,\varepsilon)$ be the optimal value of this MDP, the solution to
\begin{equation}
       V(s,\varepsilon)=\max_\pi  V_\pi(s,\varepsilon),  \quad \mbox{where} \;  V_\pi(s,\varepsilon) = E_\pi \left\{\sum_{t=0}^\infty \beta^t[r(s_t,a_t)+\varepsilon_t(a_t)]|s_0=s,\varepsilon_0=\varepsilon\right\}.
\label{eq:smoothV}
\end{equation}
Under assumptions AS, CI and EV, we can write the Bellman equation (\ref{eq:bellman}) for $V(s,\varepsilon)$ and show it 
has the representation
\begin{equation} V(s,\varepsilon)=\max_{a \in A(s)}[Q(s,a)+\epsilon(a)],
\label{eq:Vrep}
\end{equation}
where $Q(s,a)$ is given by
\begin{equation}
    Q(s,a) = r(s,a) + \beta \sum_{s'} \sigma \log\left(\sum_{a' \in A(s')} \exp\{Q(s',a')/\sigma\}\right)p(s'|s,a).
\label{eq:smoothQ}
\end{equation}
The derivation of (\ref{eq:Vrep}) and (\ref{eq:smoothQ}) makes use of an important
property of the extreme value distribution for $\varepsilon$, namely that the expectation
of the expectation of $V(s,\varepsilon)$ with respect to $\varepsilon$ 
has an  analytical expression given by the so-called log-sum or ``smoothed max function'' $V(s)$ 
\begin{eqnarray}
    V(s) &\equiv &  E\left\{V(\tilde s,\tilde \varepsilon)\left|\tilde s=s\right\}\right. 
    =  \int_\varepsilon \max_{a \in A(s)} \left[Q(s,a)+\varepsilon(a)\right]F(d\varepsilon|s) \nonumber \\ 
    &=&  \sigma \log\left(\sum_{a\in A(s)} \exp\{Q(s,a)/\sigma\}\right).
\label{eq:smoothV}
\end{eqnarray}
Using equations (\ref{eq:smoothQ}) and (\ref{eq:smoothV}) we can derive a ``smoothed Bellman equation'' for $V$
\begin{equation}
    V(s)= \sigma \log\left(\sum_{a\in A(s)} \exp\left\{[r(s,a)+\beta \sum_{s'} V(s')p(s'|s,a)]/\sigma\right\}\right),
 \label{eq:smoothBellman}
\end{equation}
Equation (\ref{eq:smoothQ}) defines a ``smoothed $Q$'' as the fixed point of the smoothed contraction mapping 
$Q_\sigma=\Lambda_\sigma(Q_\sigma)$, and (\ref{eq:smoothBellman}) defines $V$ as a fixed point of the smoothed 
Bellman operator (also a contraction mapping)
$V_\sigma=\Gamma_\sigma(V_\sigma)$. These operators are equivalent  to the corresponding operators for MDP without the state variable
$\varepsilon$ in equations (\ref{eq:q1}) and (\ref{eq:q2}) and it is not hard to show that $Q_\sigma$ and $V_\sigma$ converge
to $Q$ and $V$ as $\sigma \downarrow 0$. 


The representation of $V(x,\varepsilon)$ in equation (\ref{eq:Vrep}) means that the decision rule $\pi(a|x,\varepsilon)$ 
is isomorphic to a static discrete choice model where $Q(s,a)+\varepsilon(a)$ is the utility/payoff from action $a$.
\citet{mcfadden:1973} showed that the CCP $\pi(a|s)$ is the conditional probability that the DM chooses action $a \in A(s)$ 
given the observed state $s$ and has the classic multinomial logit formula  given by
\begin{eqnarray}
     \pi(a|s)  & =& \int I\left\{Q(s,a)+\varepsilon(a) \ge \max_{a' \in A(s)}[Q(s,a')+\varepsilon(a')]\right\}F(d\varepsilon|s)\nonumber \\
 &=&{ \exp\{Q(s,a)/\sigma\} \over \sum_{a' \in A(s)} \exp\{Q(s,a')/\sigma\}}.
\label{eq:ccp}
\end{eqnarray}
It is easy to show that as $\sigma \downarrow 0$ $\pi(a|s)$ converges to a degenerate probability that equals 1 if 
$Q(s,a)$ is the largest $Q$ value  and 0 otherwise, i.e. to the optimal pure strategy $\pi$
 in equation (\ref{eq:bellman_policy}) for MDPs without the  $\varepsilon$ variable.
If follows that the choice-specific values $Q(s,a)$ in (\ref{eq:smoothQ}) and CCPs $\pi(a|s)$ in (\ref{eq:ccp})
are exactly the same as the ``soft Q'' values and optimal policy derived in the RL
literature under the assumption that agents care about entropy of the policy $\pi$, given in equation
(\ref{eq:vdefentropy}) of section 2.4.

Equation (\ref{eq:ccp}) can be considered as the analog of the {\it policy improvement\/} step for MDPs
in equation (\ref{eq:bellman_policy}). Is there an analogous equation for the policy valuation 
step (\ref{eq:policy_valuation}) for regular MDPs where the value for any policy
$\pi$ is the solution to a linear system? \citet{AguirregabiriaMira:2002} showed that the value $V_\pi$ 
corresponding to any CCP $\pi$ is given by
\begin{equation}
   V_\pi(s) =\sum_{a \in A(s)} \pi(a|s)[r(s,a)+E\{\varepsilon(a)|s,a\}+\beta \sum_{s'} V_\pi(s')p(s'|s,a)],
\label{eq:smoothpolicyvaluation}
\end{equation}
where $E\{\varepsilon(a)|s,a\}$ is the conditional expectation of $\varepsilon(a)$ given that action $a$ is chosen
in state $s$  under decision rule $\pi$. Under the extreme value assumption for $\varepsilon$, we have
$E\{\varepsilon(a)|s,a\}=-\sigma \log(\pi(a|s))$. Using this, it is easy to see that the right side hand of 
(\ref{eq:smoothpolicyvaluation}) equals the smoothed max function defining the smoothed Bellman operator in equation
(\ref{eq:smoothBellman}). Thus, we can solve the smoothed Bellman equation $V_\sigma=\Gamma_\sigma(V_\sigma)$ using 
policy iteration just as in the case of MDPs, where we cycle between policy valuation steps (\ref{eq:smoothpolicyvaluation})
and policy improvement steps (\ref{eq:ccp}) using the $Q_\pi$ implied by $V_\pi$ in equation (\ref{eq:smoothQ}).
In fact, just as in the case of MDPs, policy iteration is equivalent to Newton's method for solving the
system $V_\sigma-\Gamma_\sigma(V_\sigma)=0$. 


Equation (\ref{eq:ccp}) implies that $\sigma\log(\pi(a|s))=Q(s,a)-V(s)\equiv \mathcal{A}(s,a)$, where
$\mathcal{A}(s,a)$ is called the {\it advantage function\/} in RL. Using $\mathcal{A}(s,a)$ we can rewrite the nonlinear fixed point
equation (\ref{eq:smoothQ}) for $Q$  as a linear system depending  on $r(s,a)$ and $\log(\pi(a|s))$,
\begin{equation}
    Q(s,a) = r(s,a) + \beta \sum_{s'} \sum_{a' \in A(s')} \left[ -\sigma\log(\pi(a'|s'))+Q(s',a')\right]\pi(a'|s')p(s'|s,a).
\label{eq:linq}
\end{equation}
Equation (\ref{eq:linq}) implies that we can decompose $Q$ as $Q(s,a)=Q_r(s,a)+Q_\varepsilon(s,a)$ where $Q_r(s,a)$ is the component 
of discounted value from observed rewards $r$ and $Q_\varepsilon(s,a)$ is the component from unobserved rewards $\varepsilon$.
Further substituting expression (\ref{eq:smoothpolicyvaluation}) in place of (\ref{eq:smoothBellman}) in the equation for
$Q$ in (\ref{eq:smoothQ}), we can show that $Q_r$ and $Q_\varepsilon$ are each the solution to the following linear systems
\begin{eqnarray}
     Q_r(s,a)&=&r(s,a)+\beta \sum_{s'} \sum_{a' \in A(s')} Q_r(s',a')\pi(a'|s')p(s'|s,a) \label{eq:Qr} \\
     Q_\varepsilon(s,a)&=&\beta \sum_{s'} \sum_{a' \in A(s')}\left[-\sigma\log(\pi(a'|s'))+ Q_\varepsilon(s',a')\right]\pi(a'|s')p(s'|s,a).
\label{eq:Qeps}
\end{eqnarray}
This decomposition can be exploited to reduce the computational burden of estimating DDC models
as we will show below.

\subsection{Estimation Methods for DDCs}

The most common estimation method for DDCs is some form of  maximum likelihood.
Suppose we observe a random sample of {\it trajectories\/} i.e. data on the  consecutive states and actions taken by
 $N$  individuals $\{(s_{it},a_{it})|t=0,\ldots,T_i, \; i=1,\ldots,N\}$. Suppose  
that we smoothly parametrize $(r,p)$ using a vector unknown parameters $\theta$, which we denote by
$r_\theta$ and $p_\theta$.\footnote{\footnotesize
The discount factor  $\beta$ is often treated as known, though in 
some cases it is also estimated and thus part of $\theta$.}  
Using the implicit function theorem, we can show that $V_\theta$ and $Q_\theta$, the contraction fixed points given in equations
(\ref{eq:smoothV}) and (\ref{eq:smoothQ}), are continuously differentiable function of $\theta$.
Using these fixed points and the logit formula for $\pi(a|s)$ in equation (\ref{eq:ccp}) we can  write a full  likelihood
function for the data as a function of the unknown parameters $\theta$ given by
\begin{equation}
    L(\theta) = \prod_{i=1}^N \prod_{t=1}^{T_i}  \pi_\theta(a_{it}|s_{it})p_\theta(s_{it}|s_{it-1},a_{it-1}).
\label{eq:lf}
\end{equation}
The smoothness of $V_\theta$ and $Q_\theta$ in $\theta$ implies that the CCP $\pi_\theta$ and the
likelihood $L(\theta)$ are both smooth functions of $\theta$ so the standard asymptotic
efficiency properties of maximum likelihood apply. 

\citet{rust:1988} introduced  the Nested Fixed Point algorithm (NFXP) to compute the maximum likelihood estimator (MLE) of $\theta$. It consists of a standard outer hill climbing algorithm to maximize $L(\theta)$ over $\theta$, but at each evaluation of $\theta$ it calculates
the contraction fixed point $Q_\theta=\Lambda_\theta(Q_\theta)$ that defines the choice-specific values $Q_\theta$ where $\Lambda_\theta$
is the operator defined on the right hand side of equation (\ref{eq:smoothQ}).
Notice that $L(\theta)$ depends on $\theta$ via $Q_\theta$. \citet{sujudd:2012} proposed an alternative strategy
for computing the MLE based on the MPEC algorithm that maximizes a likelihood $L(\theta,Q)$  expressed as  a function of $(\theta,Q)$, 
subject to the fixed point constraint $Q_\theta=\Lambda_\theta(Q_\theta)$.\footnote{\footnotesize \citet{sujudd:2012} claim that the
MPEC alogithm is substantially faster than NFXP, but  \citet{sujuddcomment:2016} showed that this was due to using SA to compute
the fixed point $Q_\theta=\Lambda_\theta(Q_\theta)$ which is slow when $\beta$ is close to 1. When Newton's method is used
to compute this fixed point, the cpu times for MPEC and NFXP are roughly comparable with NFXP being faster on test problems with
large sample sizes.}

Estimation of DDC models is often done in a two step fashion, where a parametric model of the transition probability $p(s'|s,d)$
is estimated in the first stage using a partial likelihood that does not involve the CCP $\pi$ in (\ref{eq:lf}), and in  
the second stage the parameters of $r$ are estimated using another partial likelihood involving only $\pi$, treating the estimated
parameters of $p$ as if they were known. A final Newton step iteration for maximization of the full likelihood
function (\ref{eq:lf}) can be done using the estimated parameters of $p$ and $r$ from steps 1 and 2 as
starting  values to obtain fully efficient
parameter estimates with the correct asymptotic covariance matrix, i.e. it accounts for ``first stage'' estimation effort
in the parameters of $p$.

The computational burden of NFXP stems from the need to compute the inner fixed point $Q_\theta=\Lambda_\theta(Q_\theta)$ 
for each trial value of $\theta$ in the outer hill-climbing algorithm that maximizes $L(\theta)$. 
Newton's method (e.g. policy iteration) is typically used to solve for $Q_\theta$ (or alternatively for $V_\theta$ in equation
(\ref{eq:smoothpolicyvaluation})). Even though  Newton's method converges rapidly, it requires $O(|S|^3)$ operations per Newton/policy
iteration step, making NFXP burdensome for large state spaces $S$. Therefore alternative estimation methods have been proposed that 
reduce the total number of policy iteration steps required to estimate $\theta$.


\citet{hotzmiller:1993} proposed the ``CCP estimator'' 
for $r$ that uses a non-parametric estimate of $\pi(a|s)$ to avoid repeated 
solution of $Q$ over the search for $\theta$ that NFXP requires. Consider the case where $r_\theta$ is linear in parameters,
i.e. we can write $r_\theta(s,a)=r(s,a)*\theta \equiv \sum_{k=1}^K r_k(s,a)\theta_k$ for known functions
 $\vec{r}=(r_1,\ldots,r_K)$ called {\it features\/} in the IRL
literature. Then (\ref{eq:Qr}) implies that $Q_r(s,a)=R(s,a)*\theta$ where
$R$ is the solution to
\begin{equation}
   R(s,a)=\vec{r}(s,a) + \beta \sum_{s'} \sum_{a' \in A(s')} R(s',a')\pi(a'|s')p(s'|s,a) = \vec{r}(s,a)+\beta ER(s,a),
\label{eq:hequation}
\end{equation}
where $ER(s,a)$ is a shorthand for the conditional expectation of $R(s,a)$ the middle of expression in (\ref{eq:hequation}).
It follows that we only need to solve (\ref{eq:hequation}) once using a non-parametric estimate of $\hat\pi$ along with
$Q_\varepsilon$ in (\ref{eq:Qeps}) to serve as a ``correction term''  to obtain the following form for the CCP
\begin{equation}  
           \pi_\theta(a|s) = {\exp\{\hat R(s,a)*\theta+\hat Q_\varepsilon(s,a)\} \over \sum_{a' \in A(s)} 
\exp\{\hat R(s,a')*\theta+\hat Q_\varepsilon(s,a')\}},
\label{eq:ccpestimator}
\end{equation}
where $\hat R(s,a)$ is the solution to (\ref{eq:hequation}) and $\hat Q_\varepsilon$ is the solution
to (\ref{eq:Qeps}) using non-parametric first stage
estimators $\hat\pi(a|s)$ and $\hat p(s'|s,a)$. The CCP estimator $\hat\theta$ maximizes the partial likelihood $L_\pi(\theta)$ given by 
\begin{equation}
    L_\pi(\theta) = \prod_{i=1}^N \prod_{t=1}^{T_i}  \pi_\theta(a_{it}|s_{it}),
\label{eq:ccplf}
\end{equation}
where $\pi$ is given in (\ref{eq:ccpestimator}).\footnote{\footnotesize See \citet{lrspe:2022} 
for  a modified version of the CCP estimator that is ``locally robust'' i.e. it corrects for sampling
error in the first stage estimates of the functions $\hat H(s,a)$ and $Q_\varepsilon(s,a)$.}
Thus, when $r$ is a  linear combination of known features, we only  need to compute $\hat R$ and $\hat Q_\varepsilon$ once
at the start of the estimation, requiring  $K+1$ solutions of systems with $|S|$ equations and unknowns,
where $K$ is the number of features. This can be further reduced to
a single solution if we assume that $r(s,a_s)$ is known for some
``normalizing alternative'' $a_s \in A(s)$ as we discuss further below in the
section on identification of DDC models.

A cost of the CCP estimator is that it is less efficient than using NFXP to estimate the same partial likelihood criterion (\ref{eq:ccplf})
due to the noise in the first stage non-paramteric estimates $\hat\pi$ that are used to construct
the $\hat Q_\varepsilon$ function in (\ref{eq:ccpestimator}).  \citet{AguirregabiriaMira:2002} (AM) addressed this problem
by proposing an iterative estimator called {\it nested pseudo likelihood\/} (NPL)  that uses an initial non-parametric estimate $\hat\pi$ 
to compute the value function $V_{\hat\pi}$ implied by this initial estimate by solving for it in the policy valuation step
in equation (\ref{eq:smoothpolicyvaluation}), which we denote as $\hat V_{\theta,\hat\pi}$ since it depends
both on $\hat\pi$ and the parameters $\theta$ entering
the reward function. AM then  estimate  $\hat\theta$ using a pseudo-likelihood similar to (\ref{eq:ccplf})
except that the CCP $\pi_\theta(a|s)$ is evaluated using (\ref{eq:ccp}) using estimated $Q$ functions given by
\begin{equation}
                  \hat Q_\theta(s,a)=r_\theta(s,a)+\beta \sum_{s'}\hat V_{\theta,\hat\pi}(s'|s,a)p(s'|s,a),
\label{eq:qhat}
\end{equation}
Thus, the likelihood $L_\pi(\theta)$ as depends implicitly on the first stage non-parametric CCP estimate $\hat\pi$, which can be regarded
as a ``nuisance parameter''. In principle, estimation noise in $\hat\pi$ will contaminate and affect the asymptotic 
covariance matrix for the parameters of interest $\theta$. However 
AM established a key {\it zero Jacobian property\/} that at the fixed point $\partial V(s)/\partial \pi=0$ 
which  implies that  the first stage nuisance parameter $\hat\pi$
is ``Neyman orthogonal'' with respect to the parameters of interest $\theta$ in the sense of \citet{lrspe:2022}. It follows 
that the asymptotic covariance matrix
for the NPL estimator of $\theta$ is unaffected by noise in the first stage $\hat\pi$ and is the same 
as $\pi$  was known. Further, 
AM showed that NPL can be iterated: using $\hat\theta$, they can compute an updated estimate of the CCPs 
$\hat\pi=\pi_{\hat\theta})$ and repeat the policy valuation step to get a new set of $Q$ functions (\ref{eq:qhat}), and use 
these in the likelihood (\ref{eq:ccplf}) to get an updated estimate $\hat\theta$. They refer to NPL as ``swapping''
the policy iteration and maximization steps of the NFXP estimator (which does maximization in the ``outside'' loop
and policy iteration in the ``inside'' or nested fixed point loop). Swapping the order of maximization
and policy iteration reduces the total number of policy iteration steps required to 
estimate $\theta$. AM showed that if this iterative version of  NPL converges, it produces the
same  estimate $\hat\theta$ that  NFXP produces from maximization of the partial likelihood  (\ref{eq:ccplf}).
Thus the NPL estimator is as efficient as NFXP but at lower computational cost, and it is more efficient than the CCP
estimator but at a higher computational cost.\footnote{\footnotesize
See also \citet{dearingblevins:2025} who also proposed a similar iterative
procedure called {\it efficient psuedo likelihood\/} (EPL) that is also an efficient sequential estimator that 
can estimate parameters of dynamic games.}  

Bayesian approaches have been proposed for inference in DDC models using Markov Chain Monte Carlo methods that simultaneously
simulate the likelihood as well as stochastic simulations to compute estimates of value functions that are similar to
``Monte Carlo tree search'' used by RL algorithms such as the Q-learning algorithm discussed in section 2.4, 
see \citet{ijc:2009} and \citet{norets:2009}. There are also non-Bayesian simulation estimators such as
\citet{bbl:2007} that estimate $\theta$ as the value that minimizes the degree of suboptimality 
in observed actions relative to other actions not taken  using simulation estimates of $Q$ function evaluated at the actions
the DM actually an other actions not chosen.

\subsection{Function Approximation Methods for Large State Spaces}

When the state space $S$ is very large (i.e. the state space is continuous
or the number of states $|S|$ is in the hundreds of thousands
or higher), all of the approaches discussed above that require one or more solutions 
of large linear systems with $|S|$ equations and unknowns (e.g. NFXP, CCP, NPL, etc) become computationally infeasible. 
An alternative approach is to use parametric methods to approximate the value function $V$  as $V_\phi(s)=f(s,\phi)$
where $f(s,\phi)$ is a member of a flexible family such as neural networks where $\phi$ denote the weights (parameters)
defining the network, and then recasting 
the fixed point problems (\ref{eq:smoothQ}) and (\ref{eq:smoothV}) as potentially 
more tractable nonlinear least squares problems.  For example we can approximate $V$ as $V_{\hat\phi}$ where $\hat\phi$ are
the network weights that solve the  nonlinear least squares problem of minimizing the squared Bellman residuals in 
equation (\ref{eq:nlls}). 
\begin{equation}
 \rho(\theta,\phi) = \sum_{j=1}^N  [V_\phi(s_j)-\Gamma_{\theta,\sigma}(V_\phi)(s_j)]^2.
\label{eq:nlls1}
\end{equation}
Note that the smoothed Bellman operator $\Gamma_\sigma$ defined on the right hand
side of (\ref{eq:smoothV}) depends on $r$ and $p$ which in turn are function of the model parameters $\theta$, so we write it as
$\Gamma_{\theta,\sigma}$. It follows that the minimizing parameters $\hat\phi$ are implicit functions of $\theta$. 

Let $\rho(\theta,\phi)$ denote the squared Bellman residuals defined over a grid of $N$ points in the state space.
Then a  ``nested least squares'' estimator similar to NFXP can be used estimate $\theta$ in large state
spaces, where we maximize $L(\theta)$, but in the inner fixed point problem has been transformed
into the problem of minimizing $\rho(\theta,\phi)$ over the parameters $\phi$ 
for each value of the structural parameters $\theta$.   The CCP estimator can also be implemented in a similar fashion
where we parametrize $R$ and $Q_\varepsilon$ by parameters $\phi_R$ and $\phi_\varepsilon$ and solve two
nonlinear least squares problems to approximate the fixed points in (\ref{eq:hequation}) and (\ref{eq:Qeps}), respectively.
Note, that these would only need to solved {\it once\/} at the start of estimation, and the resulting solutions
$R_{\hat\phi_R}$ and $Q_{\hat\phi_\varepsilon}$ can then used to evaluate $\pi_\theta(a|s)$  in (\ref{eq:ccpestimator}).

\citet{LuoSang:2024} introduced a penalized likelihood estimator where
$\hat\theta$ and $\hat\phi$ are simultaneously chosen
to  maximize 
\begin{equation}
(\hat\theta,\hat\phi)=\argmax_{\theta,\phi} \log(L(\theta))-\lambda \rho(\theta,\phi),
\label{eq:penalizedllf}
\end{equation}
where $\lambda \ge 0$ is a penalization
parameter.  Note that the likelihood depends on both $\theta$ and $\phi$ since both the CCP $\pi$ and $Q$
are defined in terms of $V_\phi$ which depends on the network weights, $\phi$, in addition to the structural
parameters of interest, $\theta$. 
The consistency of the penalized likelihood estimator requires the use of a {\it sieve,\/} i.e
the dimension of $\phi$ must grow to infinity at certain rate with the sample
size to guarantee that the error in approximating the true fixed point by $V_\phi$  converges to zero.
As a result they refer to the estimator (\ref{eq:penalizedllf}) as SEES, for Sieve-based Efficient Estimator for Structural models.\footnote{\footnotesize 
See \citet{arcidiacono:2013} and \citet{kristensen:2021} for other examples of sieve-based estimators for DDC models.}

A drawback of the SEES estimator is that the $\phi$ parameters determining the value function $V_\phi$ 
become nuisance parameters, and for fixed $\theta$ the 
$\phi$ that maximize (\ref{eq:penalizedllf}) $\hat\phi$ that minimizes
the squared Bellman error $\rho(\theta,\phi)$ will be an implicit function of $\theta$.
As a result, computation of the asymptotic  covariance
matrix for $\theta$ requires inverting an analog of an information matrix for $(\theta,\phi)$ jointly which can be
challenging to do when $\phi$ is high dimensional.  \citet{Nguyen:2025} introduced an estimator called
{\it neural network efficient estimator\/} (NNES) that maximizes a penalized criterion
similar to the SEES estimator (\ref{eq:penalizedllf}) but in a sequential fashion similar to the NPL
estimator of AM, benefitting from the zero Jacobian property of the NPL estimator and
ensuring Neyman orthogonality between the parameters of interest $\theta$ and the nusiance
parameters $\phi$. NNES is initialized from an initial non-parametric estimator of the CCPs, $\hat\pi$, but
instead of solving for the implied value function $V$ exactly as a system
of linear equations (\ref{eq:smoothpolicyvaluation}), it uses a 
a neural network parameterization of the value function $V_\phi$ and finds
$\hat\phi$ that minimizes
the sum of squared Bellman residuals $\rho(\theta,\phi)$ in (\ref{eq:nlls1}). The resulting value function, which
also depends on $\theta$ (so we write $V_{\theta,\hat\phi}$) is then used to form the $\hat Q_\theta$ values per (\ref{eq:qhat})
which are in turn used to calculated the CCPs
$\pi_\theta(a|s)$ and the pseudo likelihood $L_\pi(\theta)$ following the same approach as the NPL
estimator. In essence, NNES is a convex combination of the NPL and SEES estimators: 1) like NPL
it is iterative,  starting from an initial nonparametric estimator $\hat\pi$, 2) unlike NPL but similar
to SEES, it does not solve for $V_{\hat\pi}$ in equation (\ref{eq:smoothpolicyvaluation}) exactly
via standard linear algebra, but rather minimizes the square Bellman residuals $\rho(\theta,\phi)$. 
The main benefit of the NNES estimator over SEES is that the zero Jacobian/Neyman orthogonality 
property implies that to get the covariance matrix for $\theta$  we only have to invert a matrix of the
same dimension as $\theta$ rather than take the $(\theta,\theta)$ block of the inverse of a matrix
with dimension equal to the sum of the number of parameters in $\theta$ and $\phi$. For flexible
and deep neural networks, the dimension of $\phi$ can be in the thousands, so this can be a significant advantage.



\subsection{The Identification Problem}

At its core, structural estimation of DDC models can be regarded as a process of inverting observed
behavior to uncover preferences. The observed behavior from the DDC model is captured by the CCP $\pi(a|s)$, 
and it can be viewed as the  ``reduced form'' model of the agent's behavior since 
with enough data on $(a,s)$ pairs, we can estimate $\pi$ nonparametrically without need for further
assumptions.  Econometric methods that focus on {\it summarizing\/} observed behavior without attempting
to provide a deeper explanation for that behavior are
referred to as reduced-form estimation (RFE). In the RL/ML literature it is referred to as ``behavioral cloning'' (BC) 
``imitation learning'' (IM), or ``apprenticeship learning'' (AL).
The analysis of identification treats $\pi$ as known because with infinite data we could
learn it perfectly.  The preceding section showed how to derive $\pi$ from the structure $\{\beta,r,p\}$, a process that we
can summarize by the mapping $\pi=f(\beta,r,p)$. 
The identification problem concerns the invertibility of this mapping: 
given $\pi$ is there a unique underlying structure $\{\beta,r,p\}$ that implies it? Unfortunately, without
further restrictions the answer is negative, as shown by \citet{hotzmiller:1993}, 
\citet{rust:1994} and \citet{magnacthesmar:2002}. 
To understand why,  consider the first step of the inversion process, uncovering  $Q$ from $\pi$.
Using the logit form of $\pi$ in (\ref{eq:ccp}) we have
\begin{equation}
               \log\left(\pi(a|s)/\pi(a'|s)\right) = Q(s,a)-Q(s,a'), \quad s \in S, \; a,a' \in A.
\label{eq:ccpinverse}
\end{equation}
We already see that information on $\pi$ is insufficient to recover all the $Q$ values: only the {\it differences\/} in
$Q$ are identified but  not all $|S||A|$ possible values of $Q(s,a)$. Indeed, the structure $\{\beta,r,p\}$ contains $1+|S|(|A|-1)+|S|^2|A|$  
parameters which is more than then $|S|(|A|-1)$ parameters of $\pi$, so we will generally have ``more equations than unknowns'' 
and the underlying structure will only be partially identified.

Identification requires additional restrictions on $\{\beta,r,p\}$. 
Empirical studies commonly impose the assumption that $r$ and $p$ have known parametric functional forms that depend on a relatively small number of parameters $\theta$. In some situations the reward function can be treated as
known, and this facilitates the identification of subjective beliefs $p(s'|s,a)$.
For example \citet{tennis:2025} are able to identify the subjective beliefs of professional tennis servers about how their serve
strategies affects their probability of winning a tennis game. These beliefs take the form of a transition $p(s'|s,a)$ where $s$ denotes
the point state of the game and $a$ denotes serve direction. Identification of beliefs is possible  because it is reasonable to assume
$r$ is known: i.e. the server earns a reward of 1 if they win the game and 0 otherwise, and due to reasonable parametric restrictions
on $p$.


In many applications researchers are willing to assume agents have rational expectations 
which implies that their subjective beliefs about state transitions $p(s'|s,a)$ correspond to the actual probability governing
these transitions and so $p$ can be estimated non-parametrically given sufficient data, and thus can be treated as known but
with restrictions on $r$ this assumption is testable.\footnote{\footnotesize
\citet{tennis:2025} strongly reject the hypothesis the many professional tennis servers have rational beliefs about their
own ability and those of their opponent, as embodied by $p(s'|s,a)$.}
Many studies typically assume that the discount factor $\beta$ is also known, however under the assumption of 
rational expectations and restrictions on rewards, $\beta$ can be identified, see \citet{abbringdaljord:2020}.

Assuming $\{\beta,p\}$ are known, identification of $r$ depends on the assumption that either it has a known parametric
functional form $r_\theta$ where the dimension of $\theta$ 
is less than $|S|(|A|-1)$, or that for each $s \in S$ there is
a normalizing or anchor action $a_s$ such that  $r(s,a_s)$ is known. The latter assumption amounts to $|S|$ restrictions
that imply there are only $|S|(|A|-1)$ unknown rewards to be recovered, the same number of free values of $\pi$. 
\citet{hotzmiller:1993} and \citet{kang:2025} show that in this ``exactly identified'' case,
we can calculate $Q(s,a_s)$ from knowledge of $r(s,a_s)$ and $\pi(a_s|s)$ using equation (\ref{eq:linq}) 
and then (\ref{eq:ccpinverse}) allows us to recover all remaining $Q$ values. 
Using the $Q$ values we can back out the remaining unknown rewards $r$ from the smooth Bellman equation (\ref{eq:smoothQ}).

Note that the assumption that we know $r(s,a_s)$ for some anchor action $a_s$ in each state $s \in S$ is far from
innocuous. If this assumption is wrong, even though the DDC model has enough parameters to perfectly fit $\pi$,
the estimated rewards will not equal the true rewards. In this case, counterfactual behavior predicted by the model
with an incorrectly identified reward function can differ from what would actually happen under the true reward.
Another way to say this is that in the absence of any prior 
restrictions on rewards, we can only
identify a set of ``obvservationally equivalent''  rewards of the form $r(s,a)+h(s)-\beta \sum_{s'} h(s')p(s'|s,a)$ for any function $h: S \to R$. The set of all of these reward functions 
is called the partially identified set, and it defines an inherent limit
on what can be learned about agents' rewards  and emphasizes the importance of having accurate prior knowledge of 
at least some aspects of the reward function 
to obtain credible counterfactual predictions and welfare impacts of changes in policy and the environment. 
\citet{souza:2021} provides examples of certain counterfactuals that are  identified even when $r$ is
only partially identified. They conclude that their ``results call for caution while leaving room for optimism: 
although counterfactual behavior and welfare can be sensitive to identifying restrictions imposed on the model, 
there exists important classes of counterfactuals that are robust to such restrictions.'' (p. 385).



